\section{Introduction}

Language-model systems can answer from parameters, retrieve external evidence, cache reusable facts, or update themselves temporarily or durably. Existing work usually studies retrieval decisions \citep{asai2024selfrag,yan2024crag}, editing mechanisms \citep{wang2024wise,yu2023melo}, or broad editing behavior \citep{zhang2024keStudy} in isolation. We instead ask a controller question: when should one policy choose among parametric answering, retrieval, memory reads, memory writes, temporary adaptation, and guarded durable consolidation under a single objective?

Our final answer is intentionally narrower than a broad method-win claim. The frozen official results show one clear positive regime and several clear negative ones. On PopQA recurring-memory \citep{mallen2022trust}, delayed utility is useful: the full router matches retrieval quality while cutting retrieval calls substantially, and removing delayed utility collapses the result. On public FreshQA-style changing-answer evaluation derived from weekly releases \citep{vu2024freshllms}, retrieval dominates across all four official models. On MQuAKE \citep{zhong2023mquake}, retrieval again dominates, and suppressing delayed utility improves over the full router without overtaking retrieval. On UniEdit \citep{chen2025uniedit}, almost every policy is already at ceiling, so the benchmark serves mainly as a sanity check.

These findings support a regime-boundary claim:
\begin{quote}
Delayed utility helps in stable recurring-memory settings, but hurts in public changing-answer and edit-propagation settings; in those regimes, retrieval-dominant behavior wins, and removing delayed utility is more reliable than softly calibrating it.
\end{quote}

This paper makes four contributions:
\begin{itemize}[leftmargin=1.5em]
\item a unified six-action routing formulation over parameters, retrieval, reusable memory, temporary adaptation, and guarded durable consolidation;
\item a clean positive result on stable recurring-memory PopQA, where delayed utility preserves quality while reducing repeated retrieval;
\item a boundary result showing that retrieval dominates public changing-answer and propagation-heavy regimes, with soft calibration offering only limited repair relative to outright suppression of delayed utility;
\item a frozen official evidence package spanning FreshQA public main and sequence, MQuAKE, UniEdit, a manual FreshQA audit, and benchmark reliability notes.
\end{itemize}
